\section{Uniforme Discreta}
Distribuisce uniformemente la probabilità per tutti gli $n$ elementi: hanno tutti la stessa probabilità di realizzarsi.
\[\mathcal U (a, b)\]

\begin{figure}[H]
\centering

\begin{minipage}{0.45\textwidth}
\centering
\begin{tikzpicture}
\begin{axis}[
    axis lines = middle,
    ylabel = {$f(x)$},
    xlabel = {$x$},
    width = 0.90\textwidth,
    height = 0.90\textwidth,
    ymin = 0, ymax = 0.12,
    xmin = -1, xmax = 26,
    legend style={at={(axis cs: 18,0.11)}, anchor=north west, font=\tiny},
    legend cell align=left,
]

% --- distrib 1 ---
\addplot[
    only marks,
    black,
    mark=*,
    mark options={fill=black, scale=0.7}
]
coordinates {
(0,0.05)(1,0.05)(2,0.05)(3,0.05)(4,0.05)
(5,0.05)(6,0.05)(7,0.05)(8,0.05)(9,0.05)
(10,0.05)(11,0.05)(12,0.05)(13,0.05)(14,0.05)
(15,0.05)(16,0.05)(17,0.05)(18,0.05)(19,0.05)
(20,0)(21,0)(22,0)(23,0)(24,0)(25,0)
};
%\addlegendentry{$a=0,\, b=20$}

% --- distrib 2 ---
\addplot[
    only marks,
    gray!60,
    mark=*,
    mark options={fill=gray!60, scale=0.7}
]
coordinates {
(0,0)(1,0)(2,0)(3,0)
(4,0.10)(5,0.10)(6,0.10)(7,0.10)(8,0.10)
(9,0.10)(10,0.10)(11,0.10)(12,0.10)(13,0.10)(14,0.10)
(15,0)(16,0)(17,0)(18,0)(19,0)
(20,0)(21,0)(22,0)(23,0)(24,0)(25,0)
};
%\addlegendentry{$a=5,\, b=15$}

\end{axis}
\end{tikzpicture}
\end{minipage}
\hfill
\begin{minipage}{0.45\textwidth}
\centering

\begin{tikzpicture}
\begin{axis}[
    axis lines = middle,
    ylabel = {$F(k)$},
    xlabel = {$k$},
    width = 0.90\textwidth,
    height = 0.90\textwidth,
    ymin = 0, ymax = 1.1,
    xmin = -1, xmax = 26,
    legend style={at={(axis cs: 18,0.25)}, anchor=north west, font=\tiny},
    legend cell align=left,
]

% --- F(k) dist 1 ---
\addplot[
    jump mark left,
    black,
    mark=*,
    mark options={fill=black, scale=0.7}
]
coordinates{
(0,0.05)(1,0.10)(2,0.15)(3,0.20)(4,0.25)
(5,0.30)(6,0.35)(7,0.40)(8,0.45)(9,0.50)
(10,0.55)(11,0.60)(12,0.65)(13,0.70)(14,0.75)
(15,0.80)(16,0.85)(17,0.90)(18,0.95)(19,1.00)
(20,1)(21,1)(22,1)(23,1)(24,1)(25,1)
};
\addlegendentry{$a=0,\, b=20$}

% --- F(k) dist 2 ---
\addplot[
    jump mark left,
    gray!60,
    mark=*,
    mark options={fill=gray!60, scale=0.7}
]
coordinates{
(0,0)(1,0)(2,0)(3,0)(4,0)
(5,0.10)(6,0.20)(7,0.30)(8,0.40)(9,0.50)
(10,0.60)(11,0.70)(12,0.80)(13,0.90)(14,1.00)
(15,1)(16,1)(17,1)(18,1)(19,1)
(20,1)(21,1)(22,1)(23,1)(24,1)(25,1)
};
\addlegendentry{$a=5,\, b=15$}

\end{axis}
\end{tikzpicture}

\end{minipage}

\end{figure}

\bigskip
\begin{table}[H]
\centering
\renewcommand{\arraystretch}{2}
\begin{tabular}{@{}ll@{}}
\toprule
\textbf{Supporto} 
& $\{a, a+1, \dots, b\}\subset\mathbb{Z}$ \\

\midrule

\textbf{Funzione di Densità} 
& $\mathbb{P}(X=k)=\dfrac{1}{n}$, \quad $k=a,\dots,b$ \\

\midrule

\textbf{Funzione di Ripartizione}
& $F(k) = 
\begin{cases}
0, & k < a,\\[4pt]
\dfrac{k-a+1}{n}, & a \le k \le b,\\[6pt]
1, & k > b
\end{cases}$ \\[6pt]

\midrule

\textbf{Valore atteso} 
& $\mathbb{E}[X] = \dfrac{a+b}{2}$ \\

\midrule

\textbf{Varianza}
& $Var(X) = \dfrac{n^2 - 1}{12}$ \\

\midrule

\textbf{Funzione caratteristica}
& $\varphi_X(u) = \dfrac{1}{n} \displaystyle\sum_{k=a}^b e^{iuk}$ \\

\bottomrule
\end{tabular}
\end{table}